\problemname{Köpa fika}
\noindent

Ända sedan lilla Kalle vann en gratis biltvätt på trisslott har han börjat se livet annorlunda.
Han har nämligen börjat märka hur mycket tur han har och har därmed börjat satsa på allt
möjligt.
En bråkdel av hans bedrifter är att ha vunnit ett tuggummipaket när han slog vad med en klasskamrat
att Stephen Hawking skulle vara på Epsteins lista, kommit i tid till skolan varje dag trots han tar
sista bussen och fick $1.4$ utav $2$ poäng i högskoleprovet genom att kryssa i slumpmässiga svar.

Han tänker nu använda sin ofantliga tur för att hjälpa riksdagen: han tänker nämligen lösa en instans
av det ökända knapsack-problemet. Riksdagen har en fikabudget på $C$ kronor och deras fikaleverantör
har ett utbud av $N$ olika godispåsar. Varje påse $i$ har ett mått på hur uppskattad den är,
$p_i$ och en kostnad på $c_i$ kronor. Du kan bara köpa varje godispåse en gång. 
Målet är att köpa en delmängd av alla påsar, så att totala kostnaden inte överstiger $C$, samtidigt som
totala nivån av uppskattning är så stor som möjligt. I allmänhet är detta ett väldigt svårt problem att lösa,
men han hoppas att kunna lösa det snabbt genom att nyttja sin tur. 

Hans idé är följande: först, ta alla påsarna och ordna de i en slumpmässig ordning. Sedan, ignorera
de $K$ första påsarna. Därefter kommer han att betrakta varje påse i ordning och köpa den om han har
råd, annars skippa den och gå vidare. Han fortsätter så ända tills han betraktat alla $N-K$
godispåsar. 

Om han har tur med sin valda ordning kommer något $K$ ge det optimala svaret. Kalle har efter mycket slit
permuterat alla påsarna, och han ber nu dig om hjälp med att utföra hans algoritm för varje $K=0,1,\dots,N-1$.

\section*{Indata}
Första raden innehåller två heltal $N$ och $C$ ($1 \leq N \leq 2 \cdot 10^5$, $1 \leq C \leq 10^9$), antalet godispåsar 
och riksdagens fikabudget i kronors.

Den andra raden innehåller $N$ heltal $p_1, p_2, \dots, p_N$ ($1 \leq p_i \leq 10^9$), hur mycket godispåse $i$
uppskattas av riksdagen.

Den tredje raden innehåller $N$ heltal $c_1, c_2, \dots, c_N$ ($1 \leq c_i \leq C$), hur många kronor som godispåse $i$
kostar.

Notera att permutationen av godispåsar \textit{inte} är valt uniformt slumpmässigt. Detta beror på att Kalles tur gör
det mer sannolikt att en bra permutation väljs.

\section*{Utdata}
Skriv $N$ stycken heltal separerade av mellanslag, totala uppskattningen av delmängden godispåsar som Kalles
algoritm hittar för $K=0,1,\dots,N-1$ i ordning. 

\section*{Poängsättning}
Din lösning kommer att testas på en mängd testfallsgrupper.
För att få poäng för en grupp så måste du klara alla testfall i gruppen.

\noindent
\begin{tabular}{| l | l | p{12cm} |}
  \hline
  \textbf{Grupp} & \textbf{Poäng} & \textbf{Gränser} \\ \hline
  $1$    & $9$        & $N \leq 1000$ \\ \hline
  $2$    & $12$       & $C \leq 50$ \\ \hline
  $3$    & $11$       & $c_i \leq c_{i+1}$ för alla $i=1,2,\dots,N-1$ \\ \hline
  $4$    & $17$       & $c_i$ är valt uniformt slumpmässigt \footnote{Se https://en.wikipedia.org/wiki/Discrete\_uniform\_distribution} mellan $1$ och $C$. \\ \hline
  $5$    & $51$       & Inga ytterligare begränsningar. \\ \hline
\end{tabular}


